% SPDX-License-Identifier: CC-BY-SA-4.0
% Author: Matthieu Perrin

\begingroup

\begin{exercice}[Horloges logiques de Lamport]
  On considère l'exécution suivante, dont les pas sont soit des actions internes, représentées par $\bullet$,
  soit l'émission ou la réception d'un message, représentés par les flèches.
  
  \begin{center}
    \begin{tikzpicture}
      \draw[process] (0, 3) node[left]{$p_1$} to (7, 3);
      \draw[process] (0, 2) node[left]{$p_2$} to (7, 2);
      \draw[process] (0, 1) node[left]{$p_3$} to (7, 1);
      \draw[process] (0, 0) node[left]{$p_4$} to (7, 0);
      
      \draw (3.0, 3) node{$\bullet$} node[above]{$a_1$};
      \draw (6.0, 3) node{$\bullet$} node[above]{$a_2$};
      \draw (3.5, 2) node{$\bullet$} node[above]{$a_3$};
      \draw (6.0, 1) node{$\bullet$} node[above]{$a_4$};
      \draw (1.0, 1) node{$\bullet$} node[above]{$a_5$};
      \draw (4.0, 0) node{$\bullet$} node[above]{$a_6$};
      
      \draw[message] (1, 3) node[above]      {$m_1$} -- (5.0, 1);
      \draw[message] (2, 3) node[above]      {$m_2$} -- (2.5, 2);
      \draw[message] (2, 2) node[below left] {$m_3$} -- (2.5, 1);
      \draw[message] (4, 2) node[above right]{$m_4$} -- (4.5, 1);
      \draw[message] (3, 1) node[below]      {$m_5$} -- (4.5, 3);
      \draw[message] (1, 0) node[below]      {$m_6$} -- (2.0, 1);
    \end{tikzpicture}
  \end{center}

  \begin{question}
  \item Donnez l'ordre partiel happens before sur les pas d'émission des messages
  \end{question}

  \begin{question}
  \item Donnez l'horloge de Lamport associée à chacun des pas.
  \end{question}

  \begin{question}
  \item Donnez l'ordre partiel sur les messages associé à ces horloges
  \end{question}

  \begin{question}
  \item Donnez un ordre total compatible avec ces horloges.
  \end{question}

\end{exercice}

\endgroup
\endinput
